\section{Abstract}\label{sec:abstract}

We built a discrete-event simulation of the medium-voltage switchgear
warehouse at the Schneider Electric Manisa plant and used it to compare
six storage-assignment policies on kitting lead time, queue waiting time,
reach-truck (RT) utilization, throughput, and operator walking distance.
The model was driven by four months of real SAP dispatch data (the
ZWM92 transaction log, about 41,000 kit orders from nine product families)
and calibrated against a video time-motion study of 2,319 micro-events
recorded on the F400 production line.

The warehouse currently assigns bins through a capacity-based heuristic
with no systematic regard for production-line location. Standard
frequency-of-use (ABC and Double ABC) slotting rules do not
account for which line a material kits to, so high-frequency items can
end up scattered across corridors served by different milkrun (tow-train) delivery routes.
The practical result is that operators frequently need a reach truck
to access upper-level bins for materials that belong on the kit path of
a specific line --- generating queues at the RT fleet rather than at
walking distance. We wanted simulation evidence before committing to
any physical re-slotting effort.

We implemented the model in SimPy. Order arrivals were drawn from
distributions fitted to the ZWM92 log (empirical inter-batch gaps,
empirical batch sizes, categorical line and material weights); service
times came from lognormal distributions parameterised from the F400
video study. We ran 20 paired replications under common random numbers
for each of six policies: the existing SAP placement, a capacity
heuristic, usage-based ABC, Double ABC, a travel-distance
optimizer, and our proposed Line-aware Slotting policy. Differences in
lead time and waiting time across policies were tested with one-way
ANOVA and common-random-numbers paired $t$-tests.

Line-aware Slotting reduced mean kitting lead time from 12.37 minutes
per order (the SAP baseline) to 6.91 minutes, a reduction of about 44\%
($p < 0.001$, ANOVA $F = 47.1$, $df = 5/114$). The mechanism is
straightforward: placing fast-moving materials at hand-reachable levels
in the corridor closest to their destination line removes nearly
all reach-truck calls from the kit path. RT utilization fell from 21.7\%
to 1.3\%, and operator queue-waiting time --- which accounts for most
of the lead time --- dropped from 8.45 to 4.26 minutes per order.
Walking distance increased slightly, from 89.0 to 92.4 metres per order,
because bins are no longer globally sorted by proximity to the central
kitting point.

In a kitting warehouse where operators and reach trucks share the same
aisles, slotting should target the queueing mechanism, not just travel
distance. A policy that minimizes walking without controlling where
materials sit relative to production lines can inadvertently concentrate
reach-truck demand and double the effective lead time --- as the
travel-distance optimizer in this study demonstrated. Re-slotting the
fast-mover fraction that drives the roughly 400 kit orders dispatched
daily could reduce kit preparation time by about five and a half minutes
per order under the conditions modelled here.
